\documentclass[12pt]{article}
\usepackage[margin=1in]{geometry}
\usepackage{enumitem}
\usepackage{titlesec}
\usepackage{parskip}
\usepackage{graphicx}
\usepackage{hyperref}
\usepackage{fontenc}

\title{\textbf{Svolik 2012 RG Notes}\\
\large Svolik, Milan W. \textit{The Politics of Authoritarian Rule}.\\
Cambridge: Cambridge University Press, 2012.}
\author{}
\date{}

\begin{document}

\maketitle

\section{Main Idea}



%--------------------------------------------------
\section{General Notes}
%--------------------------------------------------

\begin{itemize}[leftmargin=*, itemsep=4pt]
    \item Two central conflicts of authoritarian politics:
    \begin{enumerate}
        \item \underline{The problem of authoritarian control}: Balancing against the majority excluded from power; maintaining authority and legitimacy over the people. Think external threats
        \item \underline{The problem of authoritarian power-sharing}: Countering challenges from other authorities/institutions who have some degree of power or influence over the current political process. Think internal threats
    \end{enumerate}
    \item Because dictatorships lack independent oversight institutions, agreements between autocratic leaders cannot be enforced by law or other ``hard-institutional'' means -- rather, they must be maintained through some degree of interpersonal agreement and a subsequent equilibrium outcome based on soft institutions. \textbf{self enforcing} (A further extension of this is the role of violence in authoritarian regimes, which replaces much of the institutional process seen in democratic systems.)
    \item Four dimensions of dictatorship:
    \begin{enumerate}
        \item Military involement in politics
        \item Restrictions on political parties
        \item Concentration of power with the dictator
        \item Personalization of political interactions
    \end{enumerate}
    \item There exists inherent tension between the dictator and their allies, as the dictator could use his position to expand his authority and opportunity for rent-seeking; allies fear they could be eventually excluded from the selectorate. The balance to this is the credible threat of removal of the dictator by their allies. This yields two distinct forms of authoritarian power-sharing:
    \begin{enumerate}
        \item \underline{Contested autocracy}: The dictator's allies can pose a revolutionary threat to the dictator that is sufficiently credible in order to force concessions. (Likely characteristic of the USSR post-Stalin)
        \item \underline{Established autocrats}: The dictator's allies do not pose a credible revolutionary threat and therefore cannot force concessions. (Common in African dictatorships)
    \end{enumerate}
    \item Though difficult, it is possible to establish institutions which alleviate the power-sharing commitment problem. Such institutions must monitor compliance and allow for penalties upon noncompliance. Two types of primary institutions along this route:
    \begin{enumerate}
        \item Deliberative and decision making institutions: when discussions and deliberations occur with some degree of transparency, it is easier to maintain appropriate monitoring of the dictator's actions. Such institutions facilitate regular interaction between the dictator and his allies and accordingly yield this transparency (at least to some extent). This largely explains the presence of legitimate legislatures and legislative/executive councils in dictatorships -- they are not necessarily elected nor directly empowered, but their existence supports a power-sharing balance.
        \item Formal institutions are less ambiguous and more broadly known by those involved, allowing for the detection of any noncompliance with the power-sharing agreement
    \end{enumerate}
    \item Allies face a collective action problem in responding to noncompliance or other opportunistic affairs: if only some subset seeks to punish the ruler for noncompliance, they may easily be bested and subsequently severely punished. Rather, enforcement of noncompliance requires a broad range of involvement that may be difficult to achieve even in the face of clear reneging.
    \item Complementary relationship between repression and cooptation, rather than a distinct trade-off between the two: the goal is to coopt those individuals and institutions which may be either sympathetic to the dictator or sufficiently influential enough to warrant the cost of this coercion, while repressing those individuals and institutions who would not be susceptible to such ``buy-in'' (or would require an unjustifiable amount of effort to secure their participation)
\end{itemize}


\section{Important Specific Factual Claims}

\begin{itemize}[leftmargin=*, itemsep=4pt]
    \item Aristotle \emph{Politics} book 3: ``We have next to consider how many forms of government there are, and what they are.... Tyranny is a kind of monarchy which has in view the interest of the monarch only; oligarchy has in view the interest of the wealthy; democracy of the needy. None of them the common good of all.''
    \item Significant discussion of the definition of dictatorship; notes that it must first be evaluated as binary and then evaluated as a degree. Fundamentally implies that the legislature is not comprised of freely and fairly elected officials (and to a smaller extent, that the executive is not elected freely and fairly OR by the free \& fair legislature).
    \item A transition from contested to established autocracy is virtually equivalent to a transition from oligarchy to autocracy
\end{itemize}

\section*{Connections}

\begin{itemize}
    \item Fourth dimension is personalization of political interactions -- ties directly back to Stokes et al 2013 on clientelism; potentially ties back to Dahl's Polyarchy as a closed hegemony
    \item Definitely well-connected to Gandhi \& Przeworski (2007) in a general sense of ``why'' authoritarian institutions matter
\end{itemize}

\section{My Critiques}

\begin{itemize}
    \item
\end{itemize}


\end{document}
